\section{مدول‌های انژکتیو و پروژکتیو؛ ادامه}

% Source: sources/lectures/18 Injective and Projective Modules 2.pdf | PDF page: 1 | item: divisible abelian group, examples
% AMBIGUITY[L18-Q01]: source writes y in D while naming the group A; y in A is an explicitly provisional reading.
\begin{frame}[t]{گروه آبلی بخش‌پذیر}
  \label{lec18:divisible-definition}
  \begin{definition}[خوانش موقت]
    گروه آبلیِ $A$ را \textbf{بخش‌پذیر} \lr{(divisible)} می‌گوییم هرگاه
    \[
      \forall y\in A\quad\forall n\in\mathbb Z\setminus\{0\}
      \quad\exists x\in A:\quad y=nx.
    \]
  \end{definition}
  \begin{example}
    $(\mathbb Q,+)$ بخش‌پذیر است، ولی $(\mathbb Z,+)$ بخش‌پذیر نیست.
  \end{example}
  \begin{alertblock}{خوانش موقتِ \lr{L18-Q01}}
    در سطر تعریفِ منبع، $y\in D$ نوشته شده است؛ در این‌جا به طور موقت
    $y\in A$ خوانده‌ایم. حرفِ نام گروه برای بازبینی ثبت شده است.
  \end{alertblock}
\end{frame}

% Source: sources/lectures/18 Injective and Projective Modules 2.pdf | PDF page: 1 | item: divisible iff injective Z-module theorem
% AMBIGUITY[L18-Q01]: uses the provisional consistent group name in the preceding definition.
\begin{frame}[t]{بخش‌پذیری و انژکتیو بودن}
  \label{lec18:divisible-iff-injective}
  \begin{theorem}
    فرض کنید $A$ یک گروه آبلی باشد. در این صورت
    $A$ بخش‌پذیر است اگر و تنها اگر به عنوان $\mathbb Z$-مدول انژکتیو باشد.
  \end{theorem}
  \textbf{جهت اولِ اثبات در جزوه.}
  فرض کنید $A$ به عنوان $\mathbb Z$-مدول انژکتیو باشد.
  برای $y\in A$ و $0\ne n\in\mathbb Z$ باید $x\in A$ بیابیم که $y=nx$.
\end{frame}

% Source: sources/lectures/18 Injective and Projective Modules 2.pdf | PDF page: 1 | item: f on (n), additive map and extension by injectivity
\begin{frame}[t]{گسترش همریختی از ایده‌آلِ $(n)$}
  \label{lec18:injective-divisible-extension}
  تعریف کنید
  \[
    f:(n)\longrightarrow A,\qquad kn\longmapsto ky.
  \]
  این نگاشت یک $\mathbb Z$-همریختی است؛ جمع را حفظ می‌کند.
  از انژکتیو بودنِ $A$، همریختیِ $g$ وجود دارد که نمودار زیر جابجایی است:
  \[
    \begin{tikzcd}[ampersand replacement=\&,column sep=large,row sep=large]
      0 \arrow[r] \& (n) \arrow[r,hook] \arrow[d,"f"']
        \& \mathbb Z \arrow[dl,"g"] \\
      \& A \&
    \end{tikzcd}
    \qquad g|_{(n)}=f.
  \]
\end{frame}

% Source: sources/lectures/18 Injective and Projective Modules 2.pdf | PDF page: 1 | item: g(n)=ng(1), marked x, and source why-f-alone question
\begin{frame}[t]{یافتن عنصرِ مورد نیاز}
  \label{lec18:injective-divisible-element}
  چون $g$ گسترشِ $f$ است،
  \[
    g(n)=f(n)=y.
  \]
  از طرفی $g$ یک $\mathbb Z$-همریختی است، پس
  \[
    y=g(n)=n\,g(1).
  \]
  با انتخابِ $x=g(1)$، داریم $y=nx$؛ بنابراین $A$ بخش‌پذیر است.
  \hfill$\qed$
  \medskip

  \alert{پرسشِ خودِ جزوه: چرا این کار فقط با $f$ انجام نمی‌شد؟}
\end{frame}

% Source: sources/lectures/18 Injective and Projective Modules 2.pdf | PDF page: 2 | item: converse, ideal I=(n), f(n)=y=nx, g(k)=kx
% AMBIGUITY[L18-Q02]: source omits I=0; transcription of this argument is restricted to n != 0 and does not supply the missing case.
\begin{frame}[t]{از بخش‌پذیری به گسترش همریختی}
  \label{lec18:divisible-injective-setup}
  فرض کنید $A$ بخش‌پذیر، $I$ ایده‌آلی از $\mathbb Z$ و
  $f:I\to A$ یک $\mathbb Z$-همریختی باشد.
  \[
    \begin{tikzcd}[ampersand replacement=\&,column sep=large,row sep=large]
      0 \arrow[r] \& I \arrow[r,hook] \arrow[d,"f"'] \& \mathbb Z \\
      \& A \&
    \end{tikzcd}
  \]
  جزوه می‌نویسد $I=(n)$. در حالتِ $n\ne0$، اگر $f(n)=y$،
  از بخش‌پذیری عنصری مانند $x\in A$ هست که $y=nx$.
  اکنون تعریف می‌کنیم
  \[
    g:\mathbb Z\longrightarrow A,\qquad k\longmapsto kx.
  \]
  \alert{\lr{L18-Q02}: حالتِ $I=0$ در این برهانِ منبع تصریح نشده است.}
\end{frame}

% Source: sources/lectures/18 Injective and Projective Modules 2.pdf | PDF page: 2 | item: restriction calculation and injectivity conclusion by ideal-extension criterion
% AMBIGUITY[L18-Q02]: zero ideal case of the source proof remains explicitly deferred.
\begin{frame}[t]{بررسیِ گسترش روی ایده‌آلِ ناصفر}
  \label{lec18:divisible-injective-restriction}
  نگاشتِ $g(k)=kx$ یک $\mathbb Z$-همریختی است.
  اگر $m\in(n)$، آنگاه $m=kn$ و
  \[
    \begin{aligned}
      g(m)&=mx=knx=k(nx)\\
          &=ky=f(kn)=f(m).
    \end{aligned}
  \]
  پس $g|_{(n)}=f$ و نمودار جابجایی است.
  \begin{alertblock}{بخشِ موکول به بازبینی: \lr{L18-Q02}}
    جزوه از این محاسبه، انژکتیو بودنِ $A$ را نتیجه می‌گیرد.
    این محاسبه برای ایده‌آل ناصفر رونویسی شده است؛ تکمیلِ حالتِ ایده‌آل صفر
    و در نتیجه تکمیل این جهت از اثبات، برای بازبینی بعدی ثبت شده است.
  \end{alertblock}
\end{frame}

% Source: sources/lectures/18 Injective and Projective Modules 2.pdf | PDF page: 2 | item: two prove-that exercises on projective abelian groups and PID modules
\begin{frame}[t]{دو تمرین دربارهٔ مدول‌های پروژکتیو}
  \label{lec18:projective-free-exercises}
  ثابت کنید:
  \begin{exercise}
    یک گروه آبلی پروژکتیو است اگر و تنها اگر آزاد باشد.
  \end{exercise}
  \begin{exercise}
    اگر $R$ یک دامنهٔ ایده‌آل اصلی \lr{(PID)} و $M$ یک $R$-مدول باشد، آنگاه
    $M$ پروژکتیو است اگر و تنها اگر آزاد باشد.
  \end{exercise}
\end{frame}

% Source: sources/lectures/18 Injective and Projective Modules 2.pdf | PDF page: 3 | item: injectivity iff every short exact sequence beginning in J splits
\begin{frame}[t]{معیار شکافته شدن برای مدول انژکتیو}
  \label{lec18:injective-splitting-theorem}
  \begin{theorem}
    فرض کنید $R$ یک حلقه و $J$ یک $R$-مدول باشد.
    شرط‌های زیر معادل‌اند:
    \begin{enumerate}
      \item $J$ انژکتیو است.
      \item هر رشتهٔ دقیق کوتاه به شکل
        \[
          0\longrightarrow J\longrightarrow B\longrightarrow C\longrightarrow0
        \]
        شکافته می‌شود.
    \end{enumerate}
  \end{theorem}
\end{frame}

% Source: sources/lectures/18 Injective and Projective Modules 2.pdf | PDF page: 3 | item: forward implication, identity map and retraction g phi = 1_J
\begin{frame}[t]{انژکتیو بودن، رشته را می‌شکافد}
  \label{lec18:injective-splitting-forward}
  \textbf{اثباتِ $(1)\Rightarrow(2)$.}
  کافی است نمودار زیر را بنویسیم:
  \[
    \begin{tikzcd}[ampersand replacement=\&,column sep=large,row sep=large]
      0 \arrow[r] \& J \arrow[r,"\varphi"] \arrow[d,"1_J"']
        \& B \arrow[r] \arrow[dl,"g"] \& C \arrow[r] \& 0 \\
      \& J \& \& \&
    \end{tikzcd}
  \]
  از انژکتیو بودن، $g:B\to J$ وجود دارد به طوری که
  \[
    g\varphi=1_J.
  \]
  پس رشته شکافته می‌شود.
\end{frame}

% Source: sources/lectures/18 Injective and Projective Modules 2.pdf | PDF page: 3 | item: reverse implication initial extension problem
\begin{frame}[t]{جهت عکس: مسئلهٔ گسترش}
  \label{lec18:splitting-injective-problem}
  \textbf{اثباتِ $(2)\Rightarrow(1)$.}
  نمودار زیر را در نظر بگیرید؛ $g$ یک‌به‌یک است:
  \[
    \begin{tikzcd}[ampersand replacement=\&,column sep=large,row sep=large]
      0 \arrow[r] \& A \arrow[r,"g"] \arrow[d,"f"']
        \& B \arrow[dl,dashed,"\theta\;?"] \\
      \& J \&
    \end{tikzcd}
  \]
  می‌خواهیم همریختیِ $\theta:B\to J$ بیابیم که
  \[
    \theta g=f.
  \]
\end{frame}

% Source: sources/lectures/18 Injective and Projective Modules 2.pdf | PDF page: 3 | item: W submodule of J direct sum B and quotient M
\begin{frame}[t]{ساخت یک مدول خارج‌قسمتی}
  \label{lec18:pushout-quotient}
  ابتدا $J\oplus B$ را در نظر بگیرید و بگذارید
  \[
    W=\{(f(a),-g(a))\mid a\in A\}\subseteq J\oplus B.
  \]
  در این صورت $W$ زیرمدولی از $J\oplus B$ است.
  سپس تعریف می‌کنیم
  \[
    M=\frac{J\oplus B}{W}.
  \]
\end{frame}

% Source: sources/lectures/18 Injective and Projective Modules 2.pdf | PDF page: 4 | item: maps T and S with canonical inclusion annotations
\begin{frame}[t]{دو همریختی به مدولِ $M$}
  \label{lec18:pushout-maps}
  دو $R$-همریختیِ زیر را در نظر بگیرید:
  \[
    \begin{aligned}
      T:J&\longrightarrow M=\frac{J\oplus B}{W},\\
      j&\longmapsto(j,0)+W=i_J(j)+W,
    \end{aligned}
  \]
  \[
    \begin{aligned}
      S:B&\longrightarrow M=\frac{J\oplus B}{W},\\
      b&\longmapsto(0,b)+W=i_B(b)+W.
    \end{aligned}
  \]
  در این‌جا $i_J$ و $i_B$ تزریق‌های طبیعی به جمع مستقیم‌اند.
\end{frame}

% Source: sources/lectures/18 Injective and Projective Modules 2.pdf | PDF page: 4 | item: square with Sg=Tf
\begin{frame}[t]{نمودارِ حاصل از خارج‌قسمت}
  \label{lec18:pushout-square}
  نمودار زیر جابجایی است:
  \[
    \begin{tikzcd}[ampersand replacement=\&,column sep=large,row sep=large]
      0 \arrow[r] \& A \arrow[r,"g"] \arrow[d,"f"']
        \& B \arrow[d,"S"] \\
      \& J \arrow[r,"T"'] \& M=(J\oplus B)/W
    \end{tikzcd}
  \]
  یعنی
  \[
    Sg=Tf.
  \]
\end{frame}

% Source: sources/lectures/18 Injective and Projective Modules 2.pdf | PDF page: 4 | item: coset calculation proving Sg=Tf
\begin{frame}[t]{بررسیِ جابجاییِ نمودار}
  \label{lec18:pushout-commutativity}
  برای هر $a\in A$ داریم
  \[
    \begin{aligned}
      Sg(a)&=S(g(a))=(0,g(a))+W\\
       &=(0,g(a))+(f(a),-g(a))+W\\
       &=(f(a),0)+W\\
       &=T(f(a))=Tf(a).
    \end{aligned}
  \]
  در گام دوم از $(f(a),-g(a))\in W$ استفاده شده است.
\end{frame}

% Source: sources/lectures/18 Injective and Projective Modules 2.pdf | PDF pages: 4-5 | item: T injective via g injective
\begin{frame}[t]{یک‌به‌یک بودنِ $T$}
  \label{lec18:pushout-monomorphism}
  اگر $T(j)=0$، آنگاه
  \[
    (j,0)+W=W.
  \]
  پس عنصری مانند $a\in A$ هست که
  \[
    (j,0)=(f(a),-g(a)).
  \]
  از $-g(a)=0$ و یک‌به‌یک بودنِ $g$ نتیجه می‌شود $a=0$.
  بنابراین
  \[
    j=f(a)=0.
  \]
  پس $T$ یک‌به‌یک است.
\end{frame}

% Source: sources/lectures/18 Injective and Projective Modules 2.pdf | PDF page: 5 | item: short exact sequence with T and quotient by Im T; retraction U
\begin{frame}[t]{استفاده از فرضِ شکافته شدن}
  \label{lec18:pushout-retraction}
  اکنون رشتهٔ دقیق کوتاهِ زیر را در نظر بگیرید:
  \[
    0\longrightarrow J\xrightarrow{\ T\ }M
    \xrightarrow{\text{طبیعی}}M/\operatorname{Im}T\longrightarrow0.
  \]
  بنابر فرض، این رشته شکافته می‌شود. بنابراین همریختیِ
  \[
    U:M\longrightarrow J
  \]
  وجود دارد به طوری که
  \[
    UT=1_J.
  \]
\end{frame}

% Source: sources/lectures/18 Injective and Projective Modules 2.pdf | PDF page: 5 | item: return to square, theta=US and extension conclusion
\begin{frame}[t]{ساختِ گسترش و پایان اثبات}
  \label{lec18:pushout-extension}
  به نمودار قبلی بازمی‌گردیم:
  \[
    \begin{tikzcd}[ampersand replacement=\&,column sep=large,row sep=large]
      0 \arrow[r] \& A \arrow[r,"g"] \arrow[d,"f"']
        \& B \arrow[d,"S"] \\
      \& J \arrow[r,shift left,"T"]
        \& M \arrow[l,shift left,"U"]
    \end{tikzcd}
  \]
  از $Sg=Tf$ و $UT=1_J$ داریم
  \[
    USg=(UT)f=1_Jf=f.
  \]
  با تعریف $\theta=US:B\to J$، رابطهٔ $\theta g=f$ برقرار است.
  پس $J$ انژکتیو است.\hfill$\qed$
\end{frame}

% Source: sources/lectures/18 Injective and Projective Modules 2.pdf | PDF page: 5 | item: product of injective modules iff each factor injective
\begin{frame}[t]{حاصل‌ضرب مدول‌های انژکتیو}
  \label{lec18:product-injectives}
  \begin{remark}
    برای خانوادهٔ $\{J_\alpha\}_{\alpha\in I}$ از $R$-مدول‌ها، حاصل‌ضربِ
    \[
      \prod_{\alpha\in I}J_\alpha
    \]
    انژکتیو است اگر و تنها اگر برای هر $\alpha\in I$، مدولِ $J_\alpha$ انژکتیو باشد.
  \end{remark}
\end{frame}

% Source: sources/lectures/18 Injective and Projective Modules 2.pdf | PDF page: 6 | item: unital R semisimple iff every module semisimple iff every module projective
\begin{frame}[t]{مشخصه‌ای برای حلقه‌های نیمه‌ساده}
  \label{lec18:semisimple-equivalences}
  \begin{theorem}
    فرض کنید $R$ حلقه‌ای یکدار باشد. شرط‌های زیر معادل‌اند:
    \begin{enumerate}
      \item $R$ نیمه‌ساده است.
      \item هر $R$-مدول نیمه‌ساده است.
      \item هر $R$-مدول پروژکتیو است.
    \end{enumerate}
  \end{theorem}
  \textbf{اثبات.} هم‌ارزیِ $(1)$ و $(2)$ قبلاً ثابت شده است.
\end{frame}

% Source: sources/lectures/18 Injective and Projective Modules 2.pdf | PDF page: 6 | item: 2=>3, arbitrary P, exact sequence, complementary submodule B'
\begin{frame}[t]{از نیمه‌سادگیِ همهٔ مدول‌ها به پروژکتیو بودن}
  \label{lec18:semisimple-projective-complement}
  \textbf{اثباتِ $(2)\Rightarrow(3)$.}
  فرض کنید $P$ یک $R$-مدول دلخواه باشد. نشان می‌دهیم هر رشتهٔ دقیقِ
  \[
    0\longrightarrow A\xrightarrow{\ f\ }B
      \xrightarrow{\ g\ }P\longrightarrow0
  \]
  شکافته می‌شود.
  چون $B$ نیمه‌ساده است، $f(A)\leqslant B$ مکملی دارد. \alert{چرا؟}
  پس زیرمدولی مانند $B'\leqslant B$ وجود دارد که
  \[
    B=f(A)\oplus B'.
  \]
  بنابراین هر $b\in B$ نمایش یکتایی به شکل $b=f(a)+b'$ دارد.
\end{frame}

% Source: sources/lectures/18 Injective and Projective Modules 2.pdf | PDF page: 6 | item: h:B=>A, well-definedness prompt, hf=1_A and splitting
\begin{frame}[t]{ساختِ بازگشت به $A$}
  \label{lec18:semisimple-projective-retraction}
  تعریف کنید
  \[
    h:B\longrightarrow A,\qquad f(a)+b'\longmapsto a.
  \]
  چون $f$ یک‌به‌یک است، $h$ خوش‌تعریف است. \alert{چرا؟}
  \medskip

  برای هر $a\in A$،
  \[
    hf(a)=h(f(a))=h(f(a)+0)=a.
  \]
  پس $hf=1_A$ و رشته شکافته می‌شود. بنابراین $P$ پروژکتیو است.
  \hfill$\qed$
\end{frame}

% Source: sources/lectures/18 Injective and Projective Modules 2.pdf | PDF page: 7 | item: 3=>1, arbitrary submodule I of R, R/I projective, splitting and external decomposition
\begin{frame}[t]{از پروژکتیو بودنِ همهٔ مدول‌ها به نیمه‌سادگی}
  \label{lec18:projective-semisimple-splitting}
  \textbf{اثباتِ $(3)\Rightarrow(1)$.}
  طبق فرض، $R$ به عنوان $R$-مدول پروژکتیو است.
  اگر $I$ یک زیرمدول دلخواه از $R$ باشد، رشتهٔ
  \[
    0\longrightarrow I\lhook\joinrel\longrightarrow R
      \longrightarrow R/I\longrightarrow0
  \]
  دقیق است. چون $R/I$ پروژکتیو است، این رشته شکافته می‌شود؛ بنابراین
  \[
    R\cong I\oplus R/I.
  \]
  یک یکریختی مانند $\theta:R\longrightarrow I\oplus R/I$ در نظر بگیرید.
\end{frame}

% Source: sources/lectures/18 Injective and Projective Modules 2.pdf | PDF page: 7 | item: literal theta inverse decomposition and complement conclusion
% AMBIGUITY[L18-Q03]: source does not specify theta's compatibility with I; theta^{-1}(R/I) also suppresses the summand embedding. The complement inference is deferred rather than silently repaired.
\begin{frame}[t]{گامِ مکمل‌پذیری؛ موکول به بازبینی}
  \label{lec18:projective-semisimple-complement-pending}
  جزوه پس از انتخابِ $\theta:R\to I\oplus R/I$ می‌نویسد
  \[
    R\cong I\oplus\theta^{-1}(R/I),
  \]
  و نتیجه می‌گیرد که $R$ مکمل‌پذیر و در نتیجه نیمه‌ساده است.
  \begin{alertblock}{گامِ موکول به بازبینی: \lr{L18-Q03}}
    تزریقِ عاملِ $R/I$ به جمع مستقیم و سازگاریِ $\theta$ با زیرمدولِ
    مشخصِ $I\subseteq R$ در منبع تصریح نشده‌اند.
    گذار از یکریختیِ خارجی به مکملِ همین $I$، در این رونویسی به عنوان
    گامی کامل‌شده تأیید نشده است.
  \end{alertblock}
\end{frame}
